\scanpage{30}
\mspage{10}
induit par $f$ est \uline{à coefficients entiers indépendants de $\ell$}

Situation beaucoup moins bonne en car. $p>0$, où on ne dispose pas de ``cohomologie entière'' transcendante.
Le fait que les $H^i(X,\mathbb Z_\ell)$ sont de t.f. (résultant de ce que les $H^i(X,\mathbb Z/\ell\mathbb Z)$ le sont) n'est prouvé que si $X$ est propre ou $\dim X\leq 2$ — pour l'avoir en général, il semble qu'il faille la résolution des singularités.
D'autre part, $\overline{\mathrm m}$ si $X$ est projectif et lisse, on ignore si p.ex. le rang de $H^2(X,\mathbb Z_\ell)$ est indépendant de $\ell$, et à fortiori, si le résultat plus haut pour les coefficients des polynômes caractéristiques est valable.
Cela semble actuellement, \inserted{d'ailleurs,} le \deleted{premier} plus important pas \unsure{restant} en direction d'une solution des conjectures de Weil.
\deleted{\unsure{Pour}} On peut commencer par tensoriser par $\mathbb Q_\ell$,
\margin{29}
